\documentclass[11pt,a4paper,modern]{FFExercises}
\usepackage[english,greek]{babel}
\usepackage[utf8]{inputenc}
\usepackage{nimbusserif}
\usepackage[T1]{fontenc}
\usepackage{amsmath}
\let\myBbbk\Bbbk
\let\Bbbk\relax
\usepackage[amsbb,subscriptcorrection,zswash,mtpcal,mtphrb,mtpfrak]{mtpro2}
\usepackage{graphicx,multicol,multirow,enumitem,tabularx,mathimatika,gensymb,venndiagram,hhline,longtable,tkz-euclide,fontawesome5,eurosym,tcolorbox,tabularray,tikzpagenodes,relsize,siunitx,eurosym}
\definecolor{xrwma}{HTML}{aa1212}
\definecolor{maincolor}{HTML}{aa1212}
\usetikzlibrary{calc}
\usetikzlibrary{positioning}
\tcbuselibrary{skins,theorems,breakable}
\renewcommand{\textstigma}{\textsigma\texttau}
\renewcommand{\textdexiakeraia}{}
\usepackage{svg}
\ekthetesdeiktes
\begin{document}

\titlos{Μαθηματικά}{Γ' ΕΠΑΛ}{Λύσεις - Μέτρα Διασποράς}

\askhsh
\textbf{Λύση:}
\begin{enumerate}
    \item Σωστό (\textbf{Σ})
    \item Σωστό (\textbf{Σ})
    \item Σωστό (\textbf{Σ})
    \item Λάθος (\textbf{Λ}) ($s=0$ αν όλες οι παρατηρήσεις είναι ίσες)
    \item Σωστό (\textbf{Σ})
    \item Λάθος (\textbf{Λ}) (Η τυπική απόκλιση είναι πάντα θετική, άρα $|-2| \cdot s = 2s$)
    \item Λάθος (\textbf{Λ}) (Ομοιογενές αν $CV \le 10\%$)
    \item Σωστό (\textbf{Σ})
    \item Λάθος (\textbf{Λ}) (Η διακύμανση έχει μονάδες το τετράγωνο της μεταβλητής)
    \item Λάθος (\textbf{Λ}) (Μικρότερη $s$ σημαίνει μικρότερη διασπορά)
\end{enumerate}

\askhsh
\textbf{Λύση:}
\begin{enumerate}
    \item Σωστό (\textbf{Σ})
    \item Σωστό (\textbf{Σ})
    \item Σωστό (\textbf{Σ})
    \item Σωστό (\textbf{Σ})
    \item Σωστό (\textbf{Σ})
    \item Λάθος (\textbf{Λ}) (Το $CV$ παραμένει σταθερό στον πολλαπλασιασμό με θετικό αριθμό)
    \item Σωστό (\textbf{Σ}) (Είναι και τα δύο μη αρνητικά)
    \item Λάθος (\textbf{Λ})
    \item Σωστό (\textbf{Σ}) ($s^2 = 2^2 = 4$)
    \item Σωστό (\textbf{Σ}) ($CV = \frac{s}{\bar{x}} = \frac{s}{s} = 1 = 100\%$)
\end{enumerate}

\askhsh
\textbf{Λύση:}
\begin{enumerate}
    \item \textbf{iii (Διάμεσος)} - Είναι μέτρο θέσης.
    \item \textbf{iii} - Τη διασπορά γύρω από τη μέση τιμή.
    \item \textbf{i (5)} - $s = \sqrt{25} = 5$.
    \item \textbf{iv} - $s/|\bar{x}|$.
    \item \textbf{ii} - $CV \le 10\%$ (ή $<10\%$).
\end{enumerate}

\askhsh
\textbf{Λύση:}
\begin{enumerate}
    \item \textbf{iii} - Η τυπική απόκλιση δεν επηρεάζεται από την πρόσθεση σταθεράς.
    \item \textbf{i} - Το εύρος διπλασιάζεται ($|2| \cdot R = 2R$).
    \item \textbf{ii (8)} - $Max=9, Min=1 \Rightarrow R = 9-1=8$.
    \item \textbf{iii (10\%)} - $CV = \frac{2}{20} = 0.1 = 10\%$.
    \item \textbf{iii} - Όλες οι παρατηρήσεις είναι ίσες.
\end{enumerate}

\askhsh
\textbf{Λύση:}
\begin{enumerate}
    \item $Min=3, Max=12 \Rightarrow R = 12-3 = 9$.
    \item $Min=10, Max=10 \Rightarrow R = 0$.
    \item $Min=-5, Max=5 \Rightarrow R = 5 - (-5) = 10$.
    \item $Min=0.8, Max=4.1 \Rightarrow R = 4.1 - 0.8 = 3.3$.
\end{enumerate}

\askhsh
\textbf{Λύση:}
\begin{enumerate}
    \item $\bar{x} = \frac{4+6+8+10+12}{5} = \frac{40}{5} = 8$.
    \item $R = 12-4 = 8$.
    \item $s^2 = \frac{1}{5} [(4-8)^2 + (6-8)^2 + (8-8)^2 + (10-8)^2 + (12-8)^2]$
    $= \frac{1}{5} [16 + 4 + 0 + 4 + 16] = \frac{40}{5} = 8$.
    \item $s = \sqrt{8} \approx 2.83$.
\end{enumerate}

\askhsh
\textbf{Λύση:}
\begin{enumerate}
    \item $\bar{x} = \frac{3+7+3+7}{4} = \frac{20}{4} = 5$.
    \item $s^2 = \frac{1}{4} [(3-5)^2 + (7-5)^2 + (3-5)^2 + (7-5)^2]$
    $= \frac{1}{4} [4 + 4 + 4 + 4] = \frac{16}{4} = 4$.
    Άρα $s = \sqrt{4} = 2$.
\end{enumerate}

\askhsh
\textbf{Λύση:}
\begin{enumerate}
    \item $\bar{x} = 14$. Παρατηρούμε ότι οι τιμές ισαπέχουν ανά 2. $s^2 = \frac{16+4+0+4+16}{5} = 8$, άρα $s=\sqrt{8}$.
    \item Αν αυξηθούν κατά 2: Νέα $\bar{x}' = 14+2=16$. Νέα τυπική απόκλιση $s' = s = \sqrt{8}$ (αμετάβλητη).
\end{enumerate}

\askhsh
\textbf{Λύση:}
\begin{enumerate}
    \item $R = 25 - 10 = 15$.
    \item $\bar{x} = \frac{10+15+10+20+25+10}{6} = \frac{90}{6} = 15$.
    \item $s^2 = \frac{1}{6} [(-5)^2 + 0^2 + (-5)^2 + 5^2 + 10^2 + (-5)^2]$
    $= \frac{1}{6} [25 + 0 + 25 + 25 + 100 + 25] = \frac{200}{6} = \frac{100}{3} \approx 33.33$.
\end{enumerate}

\askhsh
\textbf{Λύση:}
\begin{enumerate}
    \item $\bar{x} = \frac{5+9+11+15}{4} = \frac{40}{4} = 10$.
    \item $\sum (x_i - \bar{x}) = (5-10) + (9-10) + (11-10) + (15-10) = -5 -1 + 1 + 5 = 0$.
    \item $s^2 = \frac{25+1+1+25}{4} = \frac{52}{4} = 13$, άρα $s=\sqrt{13}$.
\end{enumerate}

\askhsh
\textbf{Λύση:}
\begin{enumerate}
    \item $\bar{x} = \frac{\sum x_i}{\nu} = \frac{100}{10} = 10$.
    \item $s^2 = \frac{1}{10} (1090 - \frac{100^2}{10}) = \frac{1}{10} (1090 - 1000) = \frac{90}{10} = 9$.
    Ή $s^2 = \frac{1090}{10} - 10^2 = 109 - 100 = 9$.
    \item $s = \sqrt{9} = 3$.
\end{enumerate}

\askhsh
\textbf{Λύση:}
Στο δείγμα Α όλες οι τιμές είναι ίσες, άρα δεν υπάρχει διασπορά ($s_A=0$).
Στο δείγμα Β υπάρχουν διαφορές, άρα $s_B > 0$. Επομένως $s_B > s_A$.
Υπολογισμός: $s_A=0$. Για το Β: $\bar{x}=8$. $s_B^2 = \frac{4+1+0+1+4}{5} = \frac{10}{5} = 2 \Rightarrow s_B=\sqrt{2}$.

\askhsh
\textbf{Λύση:}
Γνωρίζουμε ότι $s^2 = \bar{x^2} - (\bar{x})^2$.
Δίνεται $\bar{x^2} = 50$ και $(\bar{x})^2 = 36$.
Άρα $s^2 = 50 - 36 = 14$.
Επομένως $s = \sqrt{14}$.

\askhsh
\textbf{Λύση:}
$\bar{x} = \frac{4+8+x}{3} = 6 \Leftrightarrow 12+x = 18 \Leftrightarrow x=6$.
Το δείγμα είναι $4, 6, 8$.
$s^2 = \frac{(4-6)^2 + (6-6)^2 + (8-6)^2}{3} = \frac{4+0+4}{3} = \frac{8}{3}$.
$s = \sqrt{\frac{8}{3}}$.

\askhsh
\textbf{Λύση:}
$s^2 = \frac{\sum (x_i - \bar{x})^2}{\nu} = \frac{80}{5} = 16$.
$s = \sqrt{16} = 4$.

\askhsh
\textbf{Λύση:}
Έστω $x$ ο βαθμός. Οι βαθμοί είναι 5.
$\bar{x}_{new} = \frac{12+14+15+17+x}{5} = 15 \Leftrightarrow 58+x = 75 \Leftrightarrow x=17$.
Οι βαθμοί είναι: $12, 14, 15, 17, 17$. $\bar{x}=15$.
Αποκλίσεις: $-3, -1, 0, 2, 2$. Τετράγωνα: $9, 1, 0, 4, 4$.
$s^2 = \frac{18}{5} = 3.6$. $s = \sqrt{3.6} \approx 1.9$.

\askhsh
\textbf{Λύση:}
Έστω $y_i = x_i + c$. Τότε $\bar{y} = \bar{x} + c$.
Η διακύμανση είναι $s_y^2 = \frac{1}{\nu} \sum (y_i - \bar{y})^2 = \frac{1}{\nu} \sum ((x_i+c) - (\bar{x}+c))^2 = \frac{1}{\nu} \sum (x_i - \bar{x})^2 = s_x^2$.

\askhsh
\textbf{Λύση:}
$\bar{x}=3$.
$s^2 = \frac{4+1+0+1+4}{5} = 2$.

\askhsh
\textbf{Λύση:}
Οι αριθμοί προκύπτουν από τους $1,2,3,4,5$ προσθέτοντας 100.
Όπως δείξαμε, η τυπική απόκλιση δεν αλλάζει με την πρόσθεση σταθεράς.
Άρα $s = \sqrt{2}$.

\askhsh
\textbf{Λύση:}
\begin{enumerate}
    \item $\nu = 2+5+2+1 = 10$.
    $\bar{x} = \frac{1\cdot2 + 2\cdot5 + 3\cdot2 + 4\cdot1}{10} = \frac{2+10+6+4}{10} = \frac{22}{10} = 2.2$.
    \item $s^2 = \frac{2(1-2.2)^2 + 5(2-2.2)^2 + 2(3-2.2)^2 + 1(4-2.2)^2}{10}$
    $= \frac{2(1.44) + 5(0.04) + 2(0.64) + 1(3.24)}{10}$
    $= \frac{2.88 + 0.2 + 1.28 + 3.24}{10} = \frac{7.6}{10} = 0.76$.
    \item $s = \sqrt{0.76} \approx 0.87$.
\end{enumerate}

\askhsh
\textbf{Λύση:}
\begin{enumerate}
    \item Πίνακας:
    \begin{mytblr}{
        colspec={ccccc},
        row{1}={maincolor, fg=white},
        cell{2}{1-5}={maincolor!10}
    }
    $x_i$ & 10 & 12 & 15 & 18 \\
    $\nu_i$ & 2 & 3 & 2 & 2 \\
    \end{mytblr}
    (Λείπει το 20 στο τέλος, σύνολο 10 μαθητές. Διόρθωση συχνοτήτων: 10(2), 12(3), 15(2), 18(2), 20(1)).
    \item $\bar{x} = \frac{20+36+30+36+20}{10} = \frac{142}{10} = 14.2$.
    $s^2 = \dots$.
\end{enumerate}

\askhsh
\textbf{Λύση:}
$R = 20 - 10 = 10$.
$\nu=10$. $\bar{x} = \frac{10\cdot4 + 12\cdot3 + 18\cdot2 + 20\cdot1}{10} = \frac{40+36+36+20}{10} = \frac{132}{10} = 13.2$.
$s^2 = \frac{4(10-13.2)^2 + \dots}{10} = \dots = 12.96$. $s=\sqrt{12.96}=3.6$.

\askhsh
\textbf{Λύση:}
$\nu=20$. $\sum x_i \nu_i = 0\cdot2 + 1\cdot6 + 2\cdot8 + 3\cdot4 = 0+6+16+12=34$.
$\bar{x} = \frac{34}{20} = 1.7$.
$s^2 = \frac{2(1.7)^2 + 6(0.7)^2 + 8(0.3)^2 + 4(1.3)^2}{20} = \dots = 0.81$.

\askhsh
\textbf{Λύση:}
Κάθε τιμή αυξάνεται κατά 1.
Νέα μέση τιμή $\bar{x}' = 1.7 + 1 = 2.7$.
Νέα τυπική απόκλιση $s' = s = \sqrt{0.81} = 0.9$ (αμετάβλητη).

\askhsh
\textbf{Λύση:}
\begin{enumerate}
    \item $\nu = 12 + \nu_2$. $\sum x_i \nu_i = 4 + 2\nu_2 + 18 + 8 = 30 + 2\nu_2$.
    $\bar{x} = \frac{30+2\nu_2}{12+\nu_2} = 2.5 \Rightarrow 30+2\nu_2 = 30 + 2.5\nu_2 \Rightarrow -0.5\nu_2 = 0$.
    Κάτι δεν πάει καλά στα δεδομένα ή στην πράξη.
    Διόρθωση: $30+2\nu_2 = 12 \cdot 2.5 + 2.5\nu_2 \Rightarrow 30+2\nu_2 = 30+2.5\nu_2 \Rightarrow 0.5\nu_2=0 \Rightarrow \nu_2=0$.
    (Δεκτή λύση, αν και περίεργη).
    \item Αν $\nu_2=0$, τότε $\nu=12$.
    $s^2 = \dots$.
\end{enumerate}

\askhsh
\textbf{Λύση:}
$\bar{x} = \frac{500}{50} = 10$.
$s^2 = \frac{1}{50} (5450 - \frac{500^2}{50}) = \frac{1}{50} (5450 - 5000) = \frac{450}{50} = 9$.
$CV = \frac{3}{10} = 30\%$.

\askhsh
\textbf{Λύση:}
\begin{enumerate}
    \item Η κατανομή είναι συμμετρική γύρω από το 10 (συχνότητες 5, 10, 5). Άρα $\bar{x}=10$.
    \item $s^2 = \frac{5(8-10)^2 + 10(10-10)^2 + 5(12-10)^2}{20} = \frac{5\cdot4 + 0 + 5\cdot4}{20} = \frac{40}{20} = 2$.
    $s = \sqrt{2}$.
\end{enumerate}

\askhsh
\textbf{Λύση:}
\begin{enumerate}
    \item Κέντρα: $x_1=2, x_2=6, x_3=10$.
    \item $\bar{x} = \frac{2\cdot3 + 6\cdot5 + 10\cdot2}{10} = \frac{6+30+20}{10} = 5.6$.
    \item $s^2 = \dots$.
\end{enumerate}

\askhsh
\textbf{Λύση:}
Κέντρα: $10, 20, 30$.
$\nu=20$. $\bar{x} = \frac{40+200+180}{20} = \frac{420}{20} = 21$.
$s^2 = \frac{4(11)^2 + 10(1)^2 + 6(9)^2}{20} = \frac{484 + 10 + 486}{20} = \frac{980}{20} = 49$.

\askhsh
\textbf{Λύση:}
$\nu = 50$. $\bar{x} = 40$.
$s^2 = \frac{1}{50} (85000 - \frac{(\sum x_i \nu_i)^2}{50})$. Χρειαζόμαστε το $\sum x_i \nu_i = 50 \cdot 40 = 2000$.
$s^2 = \frac{1}{50} (85000 - 80000) = \frac{5000}{50} = 100$.
$s = 10$.

\askhsh
\textbf{Λύση:}
$s^2 = \frac{90}{10} = 9 \Rightarrow s=3$.
$\bar{x} = \frac{100}{10} = 10$.
$CV = \frac{3}{10} = 30\%$.

\askhsh
\textbf{Λύση:}
Κέντρα: $15, 25, 35$. $\nu=60$.
$\bar{x} = \frac{300 + 750 + 350}{60} = \frac{1400}{60} \approx 23.33$.
Υπολογισμός $s^2$.

\askhsh
\textbf{Λύση:}
Πολλαπλασιασμός με $1.05$.
$s' = 1.05 s$.
$CV' = CV$. (Ομοιογένεια σταθερή).

\askhsh
\textbf{Λύση:}
Η Β. Η τυπική απόκλιση μετράει τη διασπορά. Όταν οι τιμές είναι απλωμένες μακριά από τη μέση τιμή (ακραίες κλάσεις), η διασπορά είναι μεγαλύτερη.

\askhsh
\textbf{Λύση:}
\begin{enumerate}
    \item $CV = \frac{s}{|\bar{x}|} = \frac{2}{40} = 0.05 = 5\%$.
    \item Επειδή $5\% < 10\%$, το δείγμα είναι ομοιογενές.
\end{enumerate}

\askhsh
\textbf{Λύση:}
\begin{enumerate}
    \item $s = \sqrt{9} = 3$.
    \item $CV = \frac{3}{15} = \frac{1}{5} = 0.2 = 20\%$. Το δείγμα δεν είναι ομοιογενές ($20\% > 10\%$).
\end{enumerate}

\askhsh
\textbf{Λύση:}
$CV = \frac{150}{1200} = \frac{15}{120} = \frac{1}{8} = 0.125 = 12.5\%$.
Επειδή $12.5\% > 10\%$, τα δεδομένα δεν είναι ομοιογενή.

\askhsh
\textbf{Λύση:}
Μαθηματικά: $CV_M = \frac{2}{14} \approx 14.3\%$.
Φυσική: $CV_\Phi = \frac{3}{16} \approx 18.75\%$.
Η Φυσική έχει μεγαλύτερη σχετική διασπορά (είναι λιγότερο ομοιογενής).

\askhsh
\textbf{Λύση:}
$CV = \frac{s}{\bar{x}} \Rightarrow 0.20 = \frac{4}{\bar{x}} \Rightarrow \bar{x} = \frac{4}{0.20} = 20$.

\askhsh
\textbf{Λύση:}
Το $CV$ δεν αλλάζει στον πολλαπλασιασμό με θετικό αριθμό, διότι πολλαπλασιάζονται και το $s$ και το $\bar{x}$ με τον ίδιο αριθμό, οπότε απλοποιείται.
Νέο $CV = 50\%$.

\askhsh
\textbf{Λύση:}
\begin{enumerate}
    \item $CV_A = \frac{4}{50} = 8\% < 10\%$ (Ομοιογενές).
    \item $CV_B = \frac{10}{80} = 12.5\% > 10\%$ (Μη ομοιογενές).
\end{enumerate}

\askhsh
\textbf{Λύση:}
Αρχικά $CV = \frac{0.8}{10} = 8\%$.
Νέα μέση τιμή $\bar{x}' = 20$. Νέα $s' = 0.8$.
Νέο $CV' = \frac{0.8}{20} = 4\%$.
Ναι, θα είναι ομοιογενές (μάλιστα περισσότερο).

\askhsh
\textbf{Λύση:}
\begin{enumerate}
    \item $\bar{x} = \frac{40}{4} = 10$.
    $s^2 = \frac{25+1+1+25}{4} = 13 \Rightarrow s \approx 3.6$.
    \item $CV = \frac{3.6}{10} = 36\% > 10\%$. Όχι ομοιογενές.
\end{enumerate}

\askhsh
\textbf{Λύση:}
Πρέπει $CV \le 10\% \Leftrightarrow \frac{3}{\bar{x}} \le 0.1 \Leftrightarrow \bar{x} \ge \frac{3}{0.1} \Leftrightarrow \bar{x} \ge 30$.

\askhsh
\textbf{Λύση:}
\begin{enumerate}
    \item Μεγαλύτερη απόλυτη διασπορά έχει το Β ($s_B=6 > s_A=2$).
    \item $CV_A = 10\%$, $CV_B = \frac{6}{40}=15\%$. Μεγαλύτερη σχετική διασπορά το Β.
\end{enumerate}

\askhsh
\textbf{Λύση:}
Μπάσκετ: $CV = \frac{10}{200} = 5\%$.
Ποδόσφαιρο: $CV = \frac{12}{180} \approx 6.6\%$.
Η ομάδα μπάσκετ παρουσιάζει μεγαλύτερη ομοιογένεια (μικρότερο CV).

\askhsh
\textbf{Λύση:}
$CV = 0.2 \Rightarrow \frac{s}{10} = 0.2 \Rightarrow s=2$.
Άρα $s^2 = 4$.

\askhsh
\textbf{Λύση:}
$s = 0.5 \bar{x}$.
$CV = \frac{0.5 \bar{x}}{\bar{x}} = 0.5 = 50\%$.
Όχι, δεν είναι ομοιογενές ($50\% > 10\%$).

\askhsh
\textbf{Λύση:}
\begin{enumerate}
    \item $CV = \frac{6}{60} = 10\%$. Είναι οριακά ομοιογενές.
    \item Αφαιρούμε 5g (σταθερά). Νέα $\bar{x}=55$. Νέα $s=6$ (αμετάβλητη).
    \item Νέο $CV = \frac{6}{55} \approx 10.9\%$. Τώρα δεν είναι ομοιογενές.
\end{enumerate}

\askhsh
\textbf{Λύση:}
$CV_1 = 2/15 \approx 13.3\%$.
$CV_2 = 3/14 \approx 21.4\%$.
Το Γ1 έχει πιο ομοιογενή βαθμολογία.

\askhsh
\textbf{Λύση:}
Όχι. Αν οι συχνότητες πολλαπλασιαστούν με $k$, οι μέσοι όροι και η διακύμανση δεν αλλάζουν (τα $k$ απλοποιούνται στον αριθμητή και παρονομαστή του τύπου). Το δείγμα είναι απλά μεγαλύτερο σε μέγεθος.

\askhsh
\textbf{Λύση:}
\begin{enumerate}
    \item Κορυφή παραβολής $x = -\frac{-4}{2} = 2$, $f(2) = 4-8+6=2$.
    \item Ρίζες: $x=2$ ή $x^2-5x+6=0 \Leftrightarrow x=2, x=3$.
    \begin{alist}
        \item Παρατηρήσεις: $2, 2, 3$.
        \item $\bar{x} = \frac{7}{3} \approx 2.33$.
        $s^2 = \frac{(2-2.33)^2 \cdot 2 + (3-2.33)^2}{3} \approx \frac{0.22 + 0.44}{3} \approx 0.22 \Rightarrow s \approx 0.47$.
        \item $CV = \frac{0.47}{2.33} \approx 20\%$. Όχι ομοιογενές.
    \end{alist}
\end{enumerate}

\askhsh
\textbf{Λύση:}
Ελέγχουμε τα CV:
$CV_A = \frac{0.1}{10} = 1\%$.
$CV_B = \frac{0.08}{5} = 1.6\%$.
Η μηχανή Α έχει μικρότερο CV, άρα είναι πιο σταθερή.

\askhsh
\textbf{Λύση:}
Όχι. Για την τυπική απόκλιση του συνόλου χρειαζόμαστε τα αθροίσματα τετραγώνων ή τις επιμέρους τυπικές αποκλίσεις και μέσες τιμές. Μόνο με του μέσους όρους δεν μπορούμε να βρούμε τη διασπορά.

\askhsh
\textbf{Λύση:}
Το σφάλμα είναι στην κατανόηση/τύπο. Η τυπική απόκλιση ισούται με την τετραγωνική ρίζα της διακύμανσης, άρα είναι \textbf{πάντα μη αρνητικός αριθμός}. Δεν γίνεται να είναι $-2$.

\askhsh
\textbf{Λύση:}
Περιμένουμε το δείγμα με το μικρότερο εύρος να έχει συγκεντρωμένες τις τιμές (μικρότερη διασπορά). Άρα το Γ. (Σημείωση: Αυτό είναι εμπειρικός κανόνας, δεν ισχύει απολύτως πάντα, αλλά είναι η αναμενόμενη απάντηση σε επίπεδο διαίσθησης).

\askhsh
\textbf{Λύση:}
$s^2=0 \Leftrightarrow s=0$, άρα όλες οι τιμές είναι ίσες.
Άρα οι αριθμοί είναι $4, 4, 4, 4, 4$.

\askhsh
\textbf{Λύση:}
Οι αριθμοί είναι $x, x+2, x+4, x+6$.
Πρόκειται για παρατηρήσεις της μορφής $y_i + x$ όπου $y_i \in \{0, 2, 4, 6\}$.
Η πρόσθεση της σταθεράς $x$ δεν επηρεάζει τη διασπορά.
$s^2 = s_y^2$, σταθερή.

\askhsh
\textbf{Λύση:}
$(10-2, 10+2) = (8, 12)$.

\askhsh
\textbf{Λύση:}
$\bar{x}' = 3\bar{x}$.
Θέλουμε $CV' = CV \Leftrightarrow \frac{s'}{\bar{x}'} = \frac{s}{\bar{x}} \Leftrightarrow \frac{s'}{3\bar{x}} = \frac{s}{\bar{x}} \Leftrightarrow s' = 3s$.
Πρέπει να τριπλασιαστεί και η τυπική απόκλιση.

\askhsh
\textbf{Λύση:}
$y_i = 2x_i + 3$.
$s_y = |2| s_x = 2 \cdot 2 = 4$.

\end{document}
